\documentclass[UTF8,12pt]{ctexart}

\usepackage{amsmath}
\usepackage{cases}
\usepackage{cite}
\usepackage{tabularx}
\usepackage{graphicx}
\usepackage[margin=1in,left=20mm,right=20mm]{geometry}
\geometry{a4paper}
\usepackage{fancyhdr}
\setlength{\headheight}{14.5pt}
\pagestyle{fancy}
\fancyhf{}


\title{实验记录}
\author{张斌}
\date{\today}
\pagenumbering{arabic}

\begin{document}

\fancyhead[L]{张斌}
\fancyhead[C]{\leftmark}
\fancyfoot[C]{\thepage}

\maketitle
\tableofcontents
\newpage

\section{氧化石墨烯的制备}

\subsection{氧化石墨烯的制备(2023/4/17)}

\begin{table}[h]
    \centering
    \begin{tabularx}{0.9\textwidth} {
        l
        | >{\raggedright\arraybackslash}X 
         }
         8:30 - 10:30 & 1g 1200目 石墨粉，1g 五氧化二磷，25 mL浓硫酸，加到250 mL烧杯中。冰浴，搅拌2小时。\\
         10:30 - 12:20 & 3g 高锰酸钾，分批加入烧杯中。\\
         12:20 - 15:20 & 加入100mL去离子水，升温到35$^\circ \mathrm{C}$，搅拌3小时。 \\
         15:20 - 15:50 & 升温至85$^\circ \mathrm{C}$，搅拌30分钟。\\
         15:50 & 冷却 \\
         16:10 & 加入适量30\%过氧化氢，直至溶液变成明黄色。
    \end{tabularx}
\end{table}

\end{document}
